% Nová příprava: 2026-10-28; 1. ročník SPŠ, etapa 13–24.
% Rozpis: 1_rocnik/tematicke_plany/Rozpis_vyuky_MAT_1_rocnik.tex, hodina 22.
% Původ: vlastní formulace výkladu a vlastní úlohy; metodická inspirace uvedenými zdroji.
% Kontrola, obsahové porovnání zdrojů a časy: planovani/DOKONCENI_PRIPRAV_13_24.md.
% Učebnice: ucebnice/03_Planimetrie/Planimetrie.tex, Trojúhelník a značení;
% ucebnice/04_Goniometrie_a_trigonometrie/Goniometrické funkce.tex,
% Tabulkové hodnoty (pouze ostré úhly) a praktické použití inverzních funkcí.
% Realisticky.cz: externi_zdroje/realisticky/01_Zakladni_poznatky/01_Opakovani/010119_Podobnosti_trojuhelniku_goniometricke_funkce/010119_Lekce.pdf
% Realisticky.cz: externi_zdroje/realisticky/01_Zakladni_poznatky/01_Opakovani/010120_Pouziti_goniometrickych_funkci_ve_fyzice/010120_Lekce.pdf
\documentclass[11pt,a4paper]{article}
% Spolecny styl pro ucitelske pripravy MAT (Work/MAT)
% Verze 1.1, 2026-10-23
\NeedsTeXFormat{LaTeX2e}
\ProvidesPackage{priprava}[2026/10/23 v1.1 Sablona ucitelskych priprav MAT]

% Zakladni typografie a cestina
\RequirePackage[T1]{fontenc}
\RequirePackage[utf8]{inputenc}
\RequirePackage[czech]{babel}
\RequirePackage{lmodern}
\RequirePackage{microtype}
\RequirePackage{amsmath,amssymb}
\RequirePackage{xcolor}
\RequirePackage{enumitem}
\RequirePackage[most]{tcolorbox}
\RequirePackage{tabularx,booktabs,array}
\RequirePackage{fancyhdr}
\RequirePackage{needspace}
\RequirePackage{geometry}
\RequirePackage[hidelinks]{hyperref}

\geometry{a4paper,top=17mm,bottom=18mm,left=23mm,right=23mm}
\definecolor{notegray}{gray}{0.39}
\setlength{\parindent}{0pt}
\setlength{\parskip}{2.5pt plus 1pt minus 0.5pt}
\setlist[itemize]{topsep=2pt,itemsep=1pt,parsep=0pt,leftmargin=1.45em}
\setlist[enumerate]{topsep=2pt,itemsep=2pt,parsep=0pt,leftmargin=1.65em}
\setcounter{secnumdepth}{0}
\renewcommand{\arraystretch}{1.12}

% Zahlavi: vlevo rocnik | globalni tema, vpravo cislo: nazev hodiny
% Zapati zustava zcela prazdne.
\pagestyle{fancy}
\fancyhf{}
\renewcommand{\headrulewidth}{0.35pt}
\renewcommand{\footrulewidth}{0pt}
\setlength{\headheight}{14pt}
\newcommand{\nastavPripravu}[4]{%
  \fancyhead[L]{\small #1\enspace\textbar\enspace #2}%
  \fancyhead[R]{\small #3: #4}%
}

% Minimalisticke cernobile boxy. Kazda definice zvlast.
\tcbset{commonlesson/.style={enhanced,breakable=false,colback=white,colframe=black,
  colbacktitle=white,coltitle=black,boxrule=0.45pt,
  arc=0.8mm,outer arc=0.8mm,left=3mm,right=3mm,top=1.6mm,bottom=1.6mm,
  toptitle=1mm,bottomtitle=0.7mm,fonttitle=\bfseries,
  attach boxed title to top left={xshift=2.5mm,yshift=-1.15mm},
  boxed title style={colback=white,colframe=white,boxrule=0pt,left=1mm,right=1mm}}}
\newtcolorbox{zadaniBox}[1][]{commonlesson,title={Zadání},#1}
\newtcolorbox{prikladBox}[1][]{commonlesson,title={Příklad},#1}
\newtcolorbox{definiceBox}[1][]{commonlesson,title={Definice},#1}

% Spolecny box pro zadani a jeho reseni, bez mezery mezi castmi.
% Pouziti:
%   \begin{prikladReseniBox}
%     Zadani ulohy.
%     \tcblower
%     \begin{align*} ... \end{align*}
%   \end{prikladReseniBox}
% Spodni cast ma jemny sedy podklad a delici linku; lze v ni sazet
% vice radku matematiky. Obe casti maji jen jeden spolecny ramecek.
\tcbset{solutionlesson/.style={
  commonlesson,
  enhanced jigsaw,
  breakable,
  colbacklower=black!3,
  segmentation style={solid,draw=black!25,line width=0.35pt},
  fontlower=\small\color{notegray},
  before lower={\noindent\textbf{Řešení:}\enspace},
  leftlower=3mm,rightlower=3mm
}}
\newtcolorbox{prikladReseniBox}[1][]{solutionlesson,title={Příklad},#1}
\newtcolorbox{zadaniReseniBox}[1][]{solutionlesson,title={Zadání},#1}

% Pomucky pro autora pripravy: sedive poznamky; \reseni pro starsi pripravy
\newcommand{\poznamka}[1]{\par\begingroup\small\color{notegray}\itshape
  \textnormal{\textbf{Poznámka pro učitele.}}\enspace #1\par\endgroup}
\newcommand{\reseni}[1]{\par\begingroup\small\color{notegray}
  \textbf{Očekávané řešení:}\enspace #1\par\endgroup}
\newcommand{\blok}[1]{\needspace{5\baselineskip}\subsection*{#1}}
\newcommand{\vyrok}[1]{\uv{#1}}

\usepackage{tikz}
\usetikzlibrary{patterns,calc,angles,quotes,babel}
\nastavPripravu{1. ročník}{Pravoúhlý trojúhelník}{22}{Pythagorova věta a úhly}
\begin{document}
\blok{Nejdříve označit pravý úhel}
U pravoúhlého trojúhelníku umíme ze dvou délek určit třetí. Důležité je
rozpoznat, která strana leží proti pravému úhlu; neřídíme se otočením obrázku.
Připomeň běžné značení: strana $a$ leží proti $A$, $b$ proti $B$, $c$ proti $C$.

\begin{definiceBox}
\textbf{Pravoúhlý trojúhelník} má jeden vnitřní úhel pravý, tedy $90^\circ$.
Zbývající dva úhly jsou ostré a jejich součet je $90^\circ$.
\end{definiceBox}
\begin{definiceBox}
\textbf{Přepona} je strana proti pravému úhlu; je nejdelší stranou pravoúhlého trojúhelníku.
\end{definiceBox}
\begin{definiceBox}
\textbf{Odvěsny} jsou dvě strany svírající pravý úhel.
\end{definiceBox}
\begin{center}
\begin{tikzpicture}[scale=.65]
\coordinate (A) at (4,0); \coordinate (B) at (0,3); \coordinate (C) at (0,0);
\draw[thick] (C)--node[below]{$b$}(A)--node[above,sloped]{$c$}(B)--node[left]{$a$}(C);
\node[below right] at (A) {$A$}; \node[above] at (B) {$B$}; \node[below left] at (C) {$C$};
\pic[draw,angle radius=3mm]{right angle=A--C--B};
\pic[draw,angle radius=7mm,"$\alpha$",angle eccentricity=1.5]{angle=B--A--C};
\pic[draw,angle radius=7mm,"$\beta$",angle eccentricity=1.5]{angle=C--B--A};
\end{tikzpicture}
\end{center}
Zde je $\gamma=90^\circ$, takže přeponou je $c$. V jinak označeném trojúhelníku
může být přeponou $a$ nebo $b$.

\textbf{Pythagorova věta:} V pravoúhlém trojúhelníku s odvěsnami $a,b$ a přeponou $c$
platí $a^2+b^2=c^2$.
Jde o součet \emph{čtverců} délek, nikoli součet délek.

\blok{Proč vztah platí: porovnání obsahů}
Čtyři shodné pravoúhlé trojúhelníky sestavíme do čtverce o straně $a+b$.
Uprostřed vznikne čtverec o straně $c$: jeho rohy jsou pravé,
protože dva ostré úhly původního trojúhelníku dávají dohromady $90^\circ$.
\begin{center}
\begin{tikzpicture}[scale=.55]
\fill[black!12] (0,0) rectangle (5,5);
\fill[white] (3,0)--(5,3)--(2,5)--(0,2)--cycle;
\draw[thick] (0,0) rectangle (5,5);
\draw[thick] (3,0)--node[right,sloped]{$c$}(5,3)--(2,5)--(0,2)--cycle;
\node at (2.5,2.5) {$c^2$};
\draw (3,0)--(3,-.12);
\node[below] at (1.5,-.1) {$a$}; \node[below] at (4,-.1) {$b$};
\node[left] at (-.15,2.5) {$a+b$};
\end{tikzpicture}
\end{center}
Celkový obsah je čtyřikrát $\frac{ab}{2}$ plus $c^2$.
Tentýž velký čtverec lze rozdělit podle délek $a,b$ na čtverce s obsahy
$a^2,b^2$ a dva obdélníky s obsahem $ab$. Proto
\[
a^2+2ab+b^2=4\cdot\frac{ab}{2}+c^2=2ab+c^2,
\qquad a^2+b^2=c^2.
\]
Roznásobovací vzorec tedy nemusíme předem znát: vysvětluje jej rozdělení plochy.

\begin{prikladReseniBox}[title={P1: určujeme přeponu},breakable=false]
Odvěsny mají délky $9\,\mathrm{cm}$ a $12\,\mathrm{cm}$. Urči přeponu.
\tcblower
$c^2=9^2+12^2=81+144=225$, tedy $c=\sqrt{225}=15\,\mathrm{cm}$.
Volíme kladnou délku. Kontrola: $15>12>9$ a $15^2=225$.
Výsledek $9+12=21$ by nevyhověl rovnici čtverců.
\end{prikladReseniBox}
\begin{prikladReseniBox}[title={P2: určujeme odvěsnu},breakable=false]
V trojúhelníku $ABC$ je $\beta=90^\circ$, $b=13\,\mathrm{cm}$ a $a=5\,\mathrm{cm}$.
Urči $c$.
\tcblower
Přepona je nyní $b$. Proto $b^2=a^2+c^2$,
$c^2=13^2-5^2=169-25=144$ a $c=12\,\mathrm{cm}$.
Kontrola $5^2+12^2=13^2$. Přičtením obou známých čtverců bychom hledali další přeponu,
nikoli chybějící odvěsnu.
\end{prikladReseniBox}
\begin{prikladReseniBox}[title={P3: úhel a technický rozměr},breakable=false]
a) Pro $\gamma=90^\circ$ a $\alpha=34^\circ$ urči $\beta$.
b) Obdélníkový plech má rozměry $450\,\mathrm{mm}$ a $600\,\mathrm{mm}$. Urči úhlopříčku.
\tcblower
a) $\beta=180^\circ-90^\circ-34^\circ=56^\circ$.
b) Úhlopříčka je přepona: $d=\sqrt{450^2+600^2}
=\sqrt{562\,500}=750\,\mathrm{mm}$. Oba vstupy jsou ve stejných jednotkách.
\end{prikladReseniBox}
\poznamka{Jádro: pojmenování stran, obsahové zdůvodnění, P1–P3 a S1 a)–c),
asi 38 min práce, organizace a závěr 7 min. S1 d), e) jsou další nabídka.
Odmocninu zde chápeme jako kladnou délku; obecná pravidla pro mocniny
a odmocniny teprve přijdou od hodiny 25. Zatím určujeme úhly ze součtu,
nikoli goniometrickými poměry z příští hodiny.}

\blok{Samostatné procvičení}
\begin{zadaniReseniBox}[title={S1: vyber správnou stranu},breakable=false]
\begin{enumerate}[label=\alph*)]
\item Urči přeponu pro odvěsny $8\,\mathrm{cm}$ a $15\,\mathrm{cm}$.
\item Urči druhou odvěsnu při přeponě $25\,\mathrm{cm}$ a odvěsně $7\,\mathrm{cm}$.
\item Je-li jeden ostrý úhel $28{,}4^\circ$, jaký je druhý?
\item Může mít pravoúhlý trojúhelník přeponu $6\,\mathrm{cm}$ a odvěsnu $8\,\mathrm{cm}$?
\item Může být trojúhelník se stranami $6,8,11\,\mathrm{cm}$ pravoúhlý?
\end{enumerate}
\tcblower
a) $\sqrt{64+225}=\sqrt{289}=17\,\mathrm{cm}$.
b) $\sqrt{625-49}=\sqrt{576}=24\,\mathrm{cm}$.
c) $90^\circ-28{,}4^\circ=61{,}6^\circ$.
d) Ne, přepona musí být nejdelší; také $6^2-8^2<0$.
e) Ne: případnou přeponou by byla nejdelší strana $11$, ale
$6^2+8^2=100\ne121=11^2$. Nutná podmínka neplatí.
\end{zadaniReseniBox}

\newpage
\blok{Volitelná rezerva}
\begin{zadaniReseniBox}[title={R1: obrácená věta a minuty},breakable=false]
\begin{enumerate}[label=\alph*)]
\item Obrácená Pythagorova věta říká: je-li v trojúhelníku pro nejdelší
stranu $c$ splněno $a^2+b^2=c^2$, je úhel proti $c$ pravý.
Použij ji na strany $7,24,25\,\mathrm{cm}$.
\item Jeden ostrý úhel má $37^\circ20'$. Urči druhý. Připomeň $1^\circ=60'$.
\end{enumerate}
\tcblower
a) $7^2+24^2=49+576=625=25^2$, trojúhelník je pravoúhlý.
Obrácenou větu lze zdůvodnit porovnáním s pravoúhlým trojúhelníkem
se stejnými odvěsnami: podle Pythagorovy věty má i stejnou třetí stranu,
proto jsou trojúhelníky shodné podle sss.
b) $90^\circ00'=89^\circ60'$, takže druhý úhel je $52^\circ40'$.
Desetinné stupně a minuty nezaměňujeme: $0{,}2^\circ=12'$, nikoli $20'$.
\end{zadaniReseniBox}
\poznamka{Rezerva 10–15 min; před testováním obrácené věty musí existovat
trojúhelník s danými kladnými délkami. Zde $7+24>25$.}

\blok{Domácí procvičení}
\begin{zadaniReseniBox}[title={D1: 12–18 minut},breakable=false]
\begin{enumerate}[label=\alph*)]
\item Urči přeponu pro odvěsny $5\,\mathrm{cm}$ a $12\,\mathrm{cm}$.
\item Urči odvěsnu, je-li přepona $10\,\mathrm{cm}$ a druhá odvěsna $6\,\mathrm{cm}$.
\item Urči úhlopříčku obdélníku $18\,\mathrm{cm}$ krát $24\,\mathrm{cm}$.
\item Druhý ostrý úhel doplň k $47^\circ$ a zvlášť k $45^\circ$.
Co lze v druhém případě říci o odvěsnách?
\item Vyvrať $c=a+b$ na odvěsnách $3\,\mathrm{cm}$ a $4\,\mathrm{cm}$.
\end{enumerate}
\tcblower
a) $\sqrt{25+144}=13\,\mathrm{cm}$.
b) $\sqrt{100-36}=8\,\mathrm{cm}$.
c) $\sqrt{324+576}=30\,\mathrm{cm}$.
d) $43^\circ$ a $45^\circ$; ve druhém případě jsou odvěsny shodné,
protože proti stejným úhlům leží stejně dlouhé strany.
e) Přepona je $5\,\mathrm{cm}$, nikoli $7\,\mathrm{cm}$;
$5^2=3^2+4^2$, zatímco $7^2\ne3^2+4^2$.
\end{zadaniReseniBox}
\end{document}
